\section{序幕：当“更清晰”遇上“无真值”}

% ---------------------------------------------------------------- 钩子
\begin{frame}{看起来更清晰，就真的看清了吗？}
  \begin{columns}[T]
    \begin{column}{0.42\textwidth}
      \small
      \textbf{工业场景}：芯片热成像检测要看清\textbf{内部热轮廓}，
      据此判读缺陷与工艺偏差。
      \medskip

      \textbf{两难}：
      \begin{itemize}
        \item LWIR 8--14\,µm 受\textbf{衍射}限制，更换高端微距镜头\textbf{增益有限、成本高昂}。
        \item 那么能否\textbf{不换硬件}：在现有 20\,µm 仪器上，用\textbf{多帧微扫描 + 算法}做 $2\times$，让轮廓更清楚？
      \end{itemize}
      \medskip

      右图：学习式重建的输出\textbf{看起来确实更锐}。
    \end{column}
    \begin{column}{0.58\textwidth}
      \centering
      \includegraphics[width=\linewidth,height=0.74\textheight,keepaspectratio]{fig05_main_visual.pdf}
    \end{column}
  \end{columns}
  \takeaway{硬件路线遇到物理与成本天花板，计算增强是更现实的方向——但“更锐”本身还不是证据。}
\end{frame}

% ---------------------------------------------------------------- 陷阱 + 研究问题
\begin{frame}{陷阱：没有真值时，“更细”可能是失效信号}
  \begin{columns}[T]
    \begin{column}{0.44\textwidth}
      \small
      \textbf{致命约束}：采集与部署的\textbf{任何阶段都没有高分辨率热成像真值}。
      \begin{itemize}
        \item 配对训练数据、全参考指标、按验证 PSNR 选模型——\textbf{全部失效}。
        \item 自然图像 SR 常以“细节更丰富”当作改进证据；
              \redflag{在无真值场景，“看起来更细”恰是幻觉的典型信号}。
      \end{itemize}
      \medskip
      右图（teaser）：随训练进行，输出细节增多，
      但不依赖真值的伪影诊断\textbf{同时持续劣化}。
    \end{column}
    \begin{column}{0.56\textwidth}
      \centering
      \includegraphics[width=\linewidth,height=0.72\textheight,keepaspectratio]{fig_nullspace_drif.pdf}
    \end{column}
  \end{columns}
  \takeaway{核心研究问题：\textbf{没有真值，如何区分“真看清”与“幻觉”？}本报告即围绕这一问题展开。}
\end{frame}
